\documentclass{article}
\usepackage{lmodern}
\usepackage{amsmath}
\usepackage{listings}
\usepackage{fullpage}
\usepackage{parskip}
\usepackage{xcolor}
\lstset{
	basicstyle=\ttfamily,
	columns=fullflexible,
	frame=single,
	breaklines=true,
	postbreak=\mbox{\textcolor{red}{$\hookrightarrow$}\space},
}
\begin{document}
\title{Experiment `p\_4\_1\_c\_powers\_of\_two\_array' Results}
\author{\today}
\date{}
\maketitle
\textbf{Experiment outcome: }SUCCESS
\\\textbf{Bad responses: }0
\\\textbf{Responses containing}~\texttt{assume}~\textbf{: }0
\\\textbf{Responses containing}~\texttt{expect}~\textbf{: }0
\\\textbf{Resolution attempts: }1
\\\textbf{Hard fails (resolution): }0
\\\textbf{Soft fails (resolution): }0
\\\textbf{Verification attempts: }1
\section*{Problem Specification}
\textbf{Problem name: }p\_4\_1\_c\_powers\_of\_two\_array
\\\textbf{Natural language statement: }Write a method with loops that computes all powers of 2 from 2\^{}0 up to 2\^{}20.
\\\textbf{Method signature: }p\_4\_1\_c\_powers\_of\_two\_array() returns (powers: array<int>)
\subsection*{Ensures}
\begin{itemize}
\item \texttt{powers.Length == 21}
\item \texttt{forall i :: 0 <= i < powers.Length ==> powers[i] == pow(2, i)}
\end{itemize}
\subsection*{Functional Code Given}
\begin{lstlisting}
function pow(base: int, exp: int): int
  requires  0 <= exp
  decreases exp
{
  if exp == 0 then 1 else base * pow(base, exp - 1)
}
\end{lstlisting}
\clearpage
\section*{GenAI interactions}
Below you will find all interactions between the `user' (program) and the `assistant' (OpenAI).
\subsection*{Program $\rightarrow$ GenAI}
\begin{lstlisting}
You are given the following task to perform in Dafny:

Write a method with loops that computes all powers of 2 from 2^0 up to 2^20.

The signature should be:

p_4_1_c_powers_of_two_array() returns (powers: array<int>)

The method should respect the following contract:

ensures powers.Length == 21, ensures forall i :: 0 <= i < powers.Length ==> powers[i] == pow(2, i)

The contract uses the following dafny code:

function pow(base: int, exp: int): int
  requires  0 <= exp
  decreases exp
{
  if exp == 0 then 1 else base * pow(base, exp - 1)
}

These function(s) must not be used in your implementation of the method.

Produce and show only the Dafny body of this method, including the curly braces that surround it. Do not show the signature nor contract. You must not use 'assume' nor 'expect' anywhere in your code.
\end{lstlisting}
\subsection*{GenAI $\rightarrow$ Program}
\textbf{System fingerprint: }null
\\\textbf{ID: }chatcmpl-Dme23bS6M65lsjFFxLdsduo5zV8eC
\\\textbf{Model: }gpt-5.5-2026-04-23
\\\textbf{Created at: }1780485999
\\\textbf{Finish reason: }stop
\\\textbf{Usage: }464 tokens in, and 3987 tokens out
\begin{lstlisting}
{
  powers := new int[21];
  var i := 0;
  var p := 1;

  while i < 21
    invariant powers.Length == 21
    invariant 0 <= i <= 21
    invariant p == pow(2, i)
    invariant forall j :: 0 <= j < i ==> powers[j] == pow(2, j)
    decreases 21 - i
  {
    var oldI := i;
    var oldP := p;
    assert oldP == pow(2, oldI);

    powers[i] := p;
    assert powers[oldI] == pow(2, oldI);

    p := 2 * p;
    i := i + 1;

    assert i == oldI + 1;
    assert i > 0;
    assert i - 1 == oldI;
    assert pow(2, i) == 2 * pow(2, i - 1);
    assert p == 2 * oldP;
    assert p == pow(2, i);

    forall j | 0 <= j < i
      ensures powers[j] == pow(2, j)
    {
      if j == oldI {
        assert powers[j] == pow(2, j);
      } else {
        assert j < oldI + 1;
        assert j != oldI;
        assert j < oldI;
        assert powers[j] == pow(2, j);
      }
    }
  }
}
\end{lstlisting}
\section*{Final Program}
\begin{lstlisting}
function pow(base: int, exp: int): int
  requires  0 <= exp
  decreases exp
{
  if exp == 0 then 1 else base * pow(base, exp - 1)
}

method p_4_1_c_powers_of_two_array() returns (powers: array<int>)
	ensures powers.Length == 21
	ensures forall i :: 0 <= i < powers.Length ==> powers[i] == pow(2, i)
{
  powers := new int[21];
  var i := 0;
  var p := 1;

  while i < 21
    invariant powers.Length == 21
    invariant 0 <= i <= 21
    invariant p == pow(2, i)
    invariant forall j :: 0 <= j < i ==> powers[j] == pow(2, j)
    decreases 21 - i
  {
    var oldI := i;
    var oldP := p;
    assert oldP == pow(2, oldI);

    powers[i] := p;
    assert powers[oldI] == pow(2, oldI);

    p := 2 * p;
    i := i + 1;

    assert i == oldI + 1;
    assert i > 0;
    assert i - 1 == oldI;
    assert pow(2, i) == 2 * pow(2, i - 1);
    assert p == 2 * oldP;
    assert p == pow(2, i);

    forall j | 0 <= j < i
      ensures powers[j] == pow(2, j)
    {
      if j == oldI {
        assert powers[j] == pow(2, j);
      } else {
        assert j < oldI + 1;
        assert j != oldI;
        assert j < oldI;
        assert powers[j] == pow(2, j);
      }
    }
  }
}
\end{lstlisting}
\section*{Total Token Usage}
\textbf{Input tokens: }464
\\\textbf{Output tokens: }3987
\\\textbf{Reasoning tokens: }3582
\\\textbf{Sum of `total tokens': }4451
\section*{Experiment Timings}
\textbf{Overall Experiment} started at 1780485999159, ended at 1780486077699, lasting 78540ms (78.54 seconds)
\\\textbf{Iteration \#1} started at 1780485999160, ended at 1780486077699, lasting 78539ms (78.54 seconds)
\end{document}
